% =====================================================================
%  VisionCart Optical — User Manual
%
%  A guide for the people who USE the shop: customers buying glasses,
%  staff serving them, and the owner running the business.
%
%  This is deliberately not a technical document. Nothing here assumes
%  any knowledge of computers beyond using a web browser.
%
%  Build:  pdflatex USER-MANUAL.tex   (run twice, for the contents page)
% =====================================================================
\documentclass[11pt,a4paper]{report}

\usepackage[T1]{fontenc}
\usepackage[utf8]{inputenc}
\usepackage{lmodern}
\usepackage[english]{babel}
\usepackage{microtype}

\usepackage[
  a4paper,
  top=2.6cm, bottom=2.6cm, left=2.7cm, right=2.7cm,
  headheight=15pt
]{geometry}

\usepackage{graphicx}
\usepackage{xcolor}
\usepackage{titlesec}
\usepackage{fancyhdr}
\usepackage{enumitem}
\usepackage{booktabs}
\usepackage{longtable}
\usepackage{array}
\usepackage{float}
\usepackage[font=small,labelfont=bf,skip=6pt]{caption}
\usepackage{tcolorbox}
\usepackage{fontawesome5}
\usepackage[hidelinks]{hyperref}

\tcbuselibrary{skins,breakable}

% ---------------------------------------------------------------------
%  Colours — taken from the shop's own palette so the manual and the
%  screens it describes look like they belong together.
% ---------------------------------------------------------------------
\definecolor{vcblue}{HTML}{0B5FA5}
\definecolor{vcdark}{HTML}{121A26}
\definecolor{vcink}{HTML}{1A2230}
\definecolor{vcgrey}{HTML}{5B6472}
\definecolor{vcrule}{HTML}{D9DEE6}
\definecolor{vcpaper}{HTML}{F4F6FA}
\definecolor{vcgreen}{HTML}{1B7A47}
\definecolor{vcamber}{HTML}{8A5A1B}
\definecolor{vcamberbg}{HTML}{FDF6E3}
\definecolor{vcgreenbg}{HTML}{E5F4EC}
\definecolor{vcbluebg}{HTML}{EAF2FA}
\definecolor{vcred}{HTML}{A32828}

\hypersetup{
  colorlinks=true,
  linkcolor=vcblue,
  urlcolor=vcblue,
  pdftitle={VisionCart Optical — User Manual},
  pdfauthor={VisionCart Optical}
}

% ---------------------------------------------------------------------
%  Headings
% ---------------------------------------------------------------------
\titleformat{\part}[display]
  {\normalfont\huge\bfseries\color{vcblue}}
  {\Large\MakeUppercase{\partname\ \thepart}}{14pt}{\Huge}
\titleformat{\chapter}[display]
  {\normalfont\LARGE\bfseries\color{vcink}}
  {\normalsize\color{vcblue}\MakeUppercase{Chapter \thechapter}}{8pt}{\LARGE}
\titlespacing*{\chapter}{0pt}{-18pt}{26pt}

\titleformat{\section}
  {\normalfont\large\bfseries\color{vcink}}{\thesection}{0.8em}{}
\titleformat{\subsection}
  {\normalfont\normalsize\bfseries\color{vcink}}{\thesubsection}{0.7em}{}

\setcounter{secnumdepth}{2}
\setcounter{tocdepth}{1}

% ---------------------------------------------------------------------
%  Running heads
% ---------------------------------------------------------------------
\pagestyle{fancy}
\fancyhf{}
\renewcommand{\headrulewidth}{0.4pt}
\renewcommand{\headrule}{\hbox to\headwidth{\color{vcrule}\leaders\hrule height \headrulewidth\hfill}}
\fancyhead[L]{\small\color{vcgrey}VisionCart Optical — User Manual}
\fancyhead[R]{\small\color{vcgrey}\leftmark}
\fancyfoot[C]{\small\color{vcgrey}\thepage}

\fancypagestyle{plain}{\fancyhf{}\renewcommand{\headrulewidth}{0pt}\fancyfoot[C]{\small\color{vcgrey}\thepage}}
\renewcommand{\chaptermark}[1]{\markboth{#1}{}}

% ---------------------------------------------------------------------
%  Callout boxes
% ---------------------------------------------------------------------
\newtcolorbox{tipbox}[1][]{
  breakable, enhanced, colback=vcbluebg, colframe=vcblue,
  boxrule=0pt, leftrule=2.5pt, arc=1pt, left=10pt, right=10pt, top=8pt, bottom=8pt,
  fonttitle=\bfseries\small, coltitle=white, title={\faLightbulb\ \ Tip}, #1}

\newtcolorbox{warnbox}[1][]{
  breakable, enhanced, colback=vcamberbg, colframe=vcamber,
  boxrule=0pt, leftrule=2.5pt, arc=1pt, left=10pt, right=10pt, top=8pt, bottom=8pt,
  fonttitle=\bfseries\small, coltitle=white, title={\faExclamationTriangle\ \ Take care}, #1}

\newtcolorbox{notebox}[1][]{
  breakable, enhanced, colback=vcgreenbg, colframe=vcgreen,
  boxrule=0pt, leftrule=2.5pt, arc=1pt, left=10pt, right=10pt, top=8pt, bottom=8pt,
  fonttitle=\bfseries\small, coltitle=white, title={\faInfoCircle\ \ Good to know}, #1}

% ---------------------------------------------------------------------
%  Shorthands
% ---------------------------------------------------------------------
\newcommand{\shot}[3]{%
  \begin{figure}[H]
    \centering
    \setlength{\fboxsep}{0pt}\setlength{\fboxrule}{0.6pt}%
    {\color{vcrule}\fbox{\includegraphics[width=#2\linewidth]{dotnet/docs/screenshots/#1}}}
    \caption{#3}
  \end{figure}}

\newcommand{\btn}[1]{\textbf{#1}}
\newcommand{\menu}[1]{\textsf{\textbf{#1}}}
\newcommand{\field}[1]{\textsf{#1}}

\setlist[itemize]{leftmargin=1.3em, itemsep=3pt, topsep=5pt}
\setlist[enumerate]{leftmargin=1.5em, itemsep=5pt, topsep=5pt}

\setlength{\parskip}{6pt}
\setlength{\parindent}{0pt}
\linespread{1.06}

% =====================================================================
\begin{document}

% ---------------------------------------------------------------------
%  Title page
% ---------------------------------------------------------------------
\begin{titlepage}
\centering
\vspace*{2.2cm}

{\color{vcblue}\Large\bfseries VISIONCART OPTICAL}\\[2.2cm]

{\Huge\bfseries User Manual}\\[10pt]
{\LARGE\color{vcgrey} How to use the shop}\\[2.6cm]

\begin{minipage}{0.78\textwidth}
\centering\large\color{vcgrey}
A plain-English guide for customers buying glasses, for the staff who
serve them, and for the owner running the business.\\[10pt]
\normalsize No technical knowledge is needed --- if you can use a web
browser, you can use everything in this book.
\end{minipage}

\vfill

{\color{vcgrey}\small
Covers the complete shop: online catalogue, virtual try-on,\\
prescriptions, appointments, orders and the back office.\\[14pt]
August 2026}

\vspace{1.4cm}
\end{titlepage}

% ---------------------------------------------------------------------
%  How to use this manual
% ---------------------------------------------------------------------
\chapter*{Before you start}
\addcontentsline{toc}{chapter}{Before you start}
\markboth{Before you start}{}

\section*{Who this manual is for}

This manual is written for three different people, and you only need to
read your own part.

\begin{longtable}{@{}p{0.28\textwidth}p{0.66\textwidth}@{}}
\toprule
\textbf{If you are\ldots} & \textbf{Read\ldots} \\
\midrule
\endhead
A customer buying glasses &
\textbf{Part I}. How to find frames, try them on with your camera, enter
your prescription and place an order. \\[6pt]
Shop staff or an optician &
\textbf{Part II}. How to handle orders, check prescriptions, run the
appointment diary and keep the catalogue up to date. \\[6pt]
The shop owner or manager &
\textbf{Part III}. Staff accounts, the record of who did what, and the
handful of things worth checking each week. \\
\bottomrule
\end{longtable}

You do not need to read the parts that are not yours.

\section*{How to read the instructions}

Anything you click or type is printed in \btn{bold}. Where a page is
shown in a picture, the picture appears directly beneath the step that
mentions it.

Three kinds of note appear throughout:

\begin{tipbox}
Something that will save you time, or a shortcut worth knowing.
\end{tipbox}

\begin{warnbox}
Something that cannot easily be undone, or a mistake that is expensive
to unpick. Read these before acting.
\end{warnbox}

\begin{notebox}
Background --- why the shop behaves the way it does. Useful, but you can
skip it and still do the job.
\end{notebox}

\section*{A word about your eyes}

VisionCart is an optical practice, not a fashion shop that happens to
sell glasses. That difference shows up everywhere in this manual:

\begin{itemize}
  \item A prescription is a \emph{medical} document. The shop treats it
        as one --- it is checked by a qualified optician before any
        lenses are made.
  \item Your prescription is never edited once it has been used on an
        order. A change creates a new version, and the old one stays on
        file exactly as it was.
  \item Nothing about your eyes ever appears in a web address, so it
        cannot leak through your browser history or be seen by anyone
        looking over your shoulder at the address bar.
\end{itemize}

\begin{warnbox}
This manual explains how to \emph{use} the shop. It is not medical
advice. If your vision has changed, book a sight test --- do not simply
reorder your old prescription.
\end{warnbox}

\clearpage
\tableofcontents
\clearpage

% =====================================================================
\part{Buying glasses}
% =====================================================================

\chapter{What VisionCart does}

VisionCart is an online optician. You can browse frames, see how they
look on your own face using your camera, enter the prescription your
optician gave you, choose your lenses and have the finished glasses
delivered.

\shot{01-home}{1.0}{The shop's front page. Everything starts from the
menu along the top: \menu{All frames}, \menu{Men}, \menu{Women},
\menu{Kids}, \menu{Virtual try-on} and \menu{Deals}.}

\section{What you can do without an account}

Almost everything. You can browse, use the virtual try-on, and place an
order as a guest without signing up.

You only need an account if you want to:

\begin{itemize}
  \item look back at your past orders,
  \item save delivery addresses so you do not retype them,
  \item book an appointment at the practice,
  \item download a copy of everything the shop holds about you.
\end{itemize}

\begin{notebox}
Even as a guest, the shop opens a \emph{patient file} for you. That is
how an optical practice keeps track of your prescriptions so your
glasses can be remade if they are damaged, and so a follow-up
appointment has your history. It is a legal requirement, not marketing.
\end{notebox}

\chapter{Finding frames you like}

Click \menu{All frames} in the menu.

\shot{02-catalogue}{1.0}{The catalogue. Filters run down the left-hand
side; the frames themselves fill the rest of the page.}

\section{Narrowing things down}

The filters on the left let you narrow the list by:

\begin{itemize}
  \item \textbf{Shape} --- rectangle, round, oval, aviator, cat-eye and so on.
  \item \textbf{Material} --- acetate, metal, titanium, and mixtures.
  \item \textbf{Rim} --- full rim, half rim, or rimless.
  \item \textbf{Who it is for} --- men's, women's, unisex or children's.
  \item \textbf{Price}.
\end{itemize}

You can combine as many filters as you like. To start again, click
\btn{Clear}.

\begin{tipbox}
If you already wear glasses you are happy with, look inside one of the
arms. There is usually a set of numbers like \texttt{52\hspace{1pt}\rule[0.55ex]{5pt}{0.5pt}\hspace{1pt}18 140}. That is the
lens width, the bridge width and the arm length in millimetres ---
filter for a similar size and the new pair will fit much like the old
one.
\end{tipbox}

\section{Looking at a single frame}

Click any frame to open its own page.

\shot{03-product}{1.0}{A frame's page. The colour swatches change the
picture; the panel on the right is where you choose lenses and enter
your prescription.}

Here you can see the different colours it comes in, its measurements,
and the price. If the frame is available for virtual try-on there is a
\btn{Try it on} button.

% ---------------------------------------------------------------------
\chapter{Trying frames on with your camera}

This is the part most people come for. The shop can show you what a
frame looks like on your own face, in real time.

\shot{04-tryon}{1.0}{The virtual try-on studio.}

\section{How to use it}

\begin{enumerate}
  \item Click \menu{Virtual try-on} in the top menu.
  \item Click \btn{Use my camera}. Your browser will ask for permission
        to use the camera --- click \btn{Allow}. Nothing happens until
        you do.
  \item Once your face appears, pick any frame from the strip along the
        bottom. It will be placed on your face at the right size.
  \item Turn your head slowly from side to side to see the frame from
        different angles.
\end{enumerate}

\begin{notebox}
\textbf{Your photograph never leaves your computer.} All the work of
finding your eyes and placing the frame happens inside your own browser.
No picture of your face is sent to the shop, stored on our servers, or
seen by anybody --- unless you deliberately press \btn{Save to my file},
in which case you are choosing to keep it with your order.
\end{notebox}

\section{If the camera will not start}

\begin{itemize}
  \item \textbf{You were not asked for permission.} You may have refused
        it before. Click the small padlock or camera icon in your
        browser's address bar and allow the camera for this site.
  \item \textbf{Another program is using the camera.} Close video-call
        software and try again.
  \item \textbf{You would rather not use the camera at all.} Use
        \btn{Upload a photo} instead. A clear, front-facing photo in
        good light works best.
\end{itemize}

\section{If the frame sits in the wrong place}

The shop finds your pupils automatically, but strong glare, a hat, or a
heavy fringe can confuse it. If the frame looks wrong:

\begin{enumerate}
  \item Click \btn{Adjust eye points}.
  \item Click once on the centre of your right pupil, then once on your
        left.
  \item The frame is repositioned around the two points you marked.
\end{enumerate}

\section{Your pupillary distance (PD)}

While the try-on is running, the shop measures the distance between the
centres of your pupils. This is your \textbf{pupillary distance}, usually
written PD, and it is measured in millimetres --- most adults are
between 54 and 74.

It matters because your lenses have to be positioned so that the optical
centre of each lens sits in front of the centre of each pupil. If the PD
is wrong, strong lenses can cause eye strain or headaches.

\begin{tipbox}
The shop shows how confident it is in the measurement. If it says the
reading is uncertain, hold your head straight, look directly at the
camera, and make sure the room is evenly lit. If your optician has
already written a PD on your prescription, use theirs --- you can type
it in by hand.
\end{tipbox}

% ---------------------------------------------------------------------
\chapter{Your prescription}

When you buy prescription lenses, you need the prescription your
optician gave you after your last sight test. It will look something
like the table below.

\begin{center}
\small
\begin{tabular}{@{}lccccc@{}}
\toprule
& \textbf{SPH} & \textbf{CYL} & \textbf{AXIS} & \textbf{ADD} & \textbf{PD} \\
\midrule
\textbf{Right (OD)} & $-2.25$ & $-0.75$ & 175 & --- & 31.5 \\
\textbf{Left (OS)}  & $-2.00$ & $-0.50$ & 10  & --- & 31.5 \\
\bottomrule
\end{tabular}
\end{center}

\section{What the words mean}

\begin{longtable}{@{}p{0.20\textwidth}p{0.74\textwidth}@{}}
\toprule
\textbf{Term} & \textbf{In plain English} \\
\midrule
\endhead
\textbf{OD} & Your \emph{right} eye. From the Latin \emph{oculus dexter}. \\[5pt]
\textbf{OS} & Your \emph{left} eye. From \emph{oculus sinister}. \\[5pt]
\textbf{SPH} (Sphere) &
How short-sighted or long-sighted you are. A minus number means you
cannot see far away clearly; a plus number means you struggle up close.
The bigger the number, the stronger the lens. \\[5pt]
\textbf{CYL} (Cylinder) &
Corrects astigmatism --- an eye shaped slightly more like a rugby ball
than a football. Many people have none, and the box is simply blank. \\[5pt]
\textbf{AXIS} &
Only used when there is a cylinder value. It is an angle between 1 and
180 that tells the lab which way round to grind the correction. \\[5pt]
\textbf{ADD} &
Extra magnifying power for reading. It appears on prescriptions for
bifocals and varifocals. \\[5pt]
\textbf{PD} &
Pupillary distance --- see the previous chapter. \\
\bottomrule
\end{longtable}

\shot{08-guide-prescription}{1.0}{The shop's own prescription guide, at
\menu{Help} $\rightarrow$ \menu{Understanding your prescription}.}

\section{Typing it in}

On the frame's page, choose \btn{Prescription lenses} and the boxes
appear. Every value is chosen from a drop-down list rather than typed
freely.

\begin{notebox}
The drop-downs only offer values that step by 0.25 --- $-2.00$,
$-2.25$, $-2.50$ and so on. This is not the shop being awkward: real
prescriptions only ever come in quarter-steps, and a typed value like
$-2.13$ is one nobody can make. Choosing from a list means a slipped
finger cannot turn into a ruined pair of glasses.
\end{notebox}

\begin{warnbox}
\textbf{Check the right and left columns carefully.} Swapping them is
the single most common mistake, and the resulting glasses will feel
wrong in a way that is hard to describe. Read your prescription once
more before you continue.
\end{warnbox}

If your prescription has a cylinder value, you must also give the axis.
The shop will not let you continue without it --- a cylinder with no
axis cannot be made.

\section{Prescriptions that have run out}

Prescriptions expire, usually after two years. If yours is out of date,
book a sight test. The shop will accept the order, but an optician will
check it before anything is made, and they may contact you.

% ---------------------------------------------------------------------
\chapter{Choosing lenses}

Once your prescription is in, you choose the lenses themselves. There
are four decisions, and the price updates as you make each one.

\section{Lens type}

\begin{itemize}
  \item \textbf{Single vision} --- one prescription across the whole
        lens. The usual choice for distance or reading glasses.
  \item \textbf{Bifocal} --- distance at the top, reading at the bottom,
        with a visible line between.
  \item \textbf{Varifocal} (progressive) --- distance, middle and
        reading, blending smoothly with no line.
\end{itemize}

\section{Thickness}

Stronger prescriptions need thicker lenses. A higher \emph{index} makes
the same prescription thinner and lighter, at a higher price.

\begin{longtable}{@{}p{0.16\textwidth}p{0.78\textwidth}@{}}
\toprule
\textbf{Index} & \textbf{Best for} \\
\midrule
\endhead
1.50 & Weak prescriptions, roughly up to $\pm 2.00$. \\
1.56 & A little thinner; a good general choice. \\
1.60 & Around $\pm 2.00$ to $\pm 4.00$. Noticeably slimmer. \\
1.67 & Around $\pm 4.00$ to $\pm 6.00$. \\
1.74 & Very strong prescriptions. The thinnest available. \\
\bottomrule
\end{longtable}

\begin{tipbox}
If your prescription is stronger than about $\pm 4.00$ and you have
chosen a large frame, it is worth paying for a higher index. A big lens
with a strong prescription gets thick at the edges, and it shows.
\end{tipbox}

\section{Coatings}

\begin{itemize}
  \item \textbf{Anti-reflective} --- cuts glare from headlights and
        screens. Recommended for almost everyone.
  \item \textbf{Scratch-resistant} --- a harder surface.
  \item \textbf{Blue-light filter} --- reduces blue light from screens.
  \item \textbf{Photochromic} --- darkens in sunlight, clears indoors.
\end{itemize}

\section{Tints}

For sunglasses: grey, brown or green, in a choice of strengths.

% ---------------------------------------------------------------------
\chapter{Placing your order}

\section{The bag}

Click \menu{Bag} in the top right to see what you have chosen.

\shot{05-cart}{1.0}{The bag. Each line shows the frame, the lenses you
chose and a summary of the prescription attached to it.}

You can change quantities or remove a line here. The prices are worked
out fresh by the shop every time you open this page.

\section{Discount codes}

If you have a code, type it into the \field{Discount code} box and click
\btn{Apply}. If it will not work, the shop says exactly why --- expired,
minimum spend not met, or not valid on those frames --- rather than
simply refusing.

\section{Checking out}

Click \btn{Checkout} and fill in:

\begin{enumerate}
  \item \textbf{Your contact details} --- email address and telephone
        number. The email address is where your confirmation and updates
        are sent, so check it carefully.
  \item \textbf{Delivery address.} If you have an account with saved
        addresses, choose one from the list.
  \item \textbf{Delivery method} --- the options and costs for your area.
  \item \textbf{How you would like to pay.}
\end{enumerate}

\section{Ways to pay}

\begin{itemize}
  \item \textbf{Cash on delivery} --- pay the courier when the glasses
        arrive.
  \item \textbf{Bank transfer} --- the shop's account details are shown
        after you order. Send your receipt quoting the order number.
  \item \textbf{Card} --- if the shop has card payments switched on.
\end{itemize}

\section{What happens next}

You will see a confirmation page with your order number, which looks
like \texttt{VC-2026-000123}, and receive an email straight away.

Then, behind the scenes:

\begin{enumerate}
  \item An optician checks your prescription. This is a real person, and
        it is the step that protects you.
  \item Once approved, the order goes to the lab and your lenses are
        made and fitted to your frame.
  \item The finished glasses are dispatched, and you get an email with
        tracking details.
\end{enumerate}

\begin{notebox}
If the optician finds a problem --- a value that looks unlikely, or a
prescription that has expired --- they will contact you rather than
guess. You are told why, and nothing is made until it is sorted out.
\end{notebox}

% ---------------------------------------------------------------------
\chapter{Your account}

\shot{10-account}{1.0}{The account page, with recent orders and links to
everything else.}

\section{Creating one}

Click \menu{Sign in} then \btn{Create one}. You need an email address
and a password of at least eight characters.

\shot{07-register}{0.92}{Creating an account.}

\section{Your orders}

Your recent orders are listed on the account page. Click any of them to
see its full detail, including where it has reached.

\section{Your addresses}

\shot{11-addresses}{1.0}{The address book. The address marked
\emph{Default} is the one checkout offers first.}

Click \btn{Add an address} to save a new one.

\shot{12-address-form}{0.92}{Adding an address.}

\begin{itemize}
  \item The \field{Label} is just for you --- ``Home'', ``Work'',
        ``Mum's''.
  \item Ticking \field{Use this as my default} makes checkout offer it
        first.
  \item \btn{Remove} takes an address out of your book. Your past orders
        keep their own record of where they were delivered, so nothing
        is lost from your order history.
\end{itemize}

\section{Changing your password}

Click \menu{Sign in}, then \btn{Forgotten your password?}. You will be
emailed a link that works once and expires after six hours.

\begin{notebox}
Changing your password signs you out everywhere else. If you have used a
shared or public computer, changing your password is the quickest way to
be certain nobody is still signed in as you.
\end{notebox}

% ---------------------------------------------------------------------
\chapter{Booking an appointment}

If the practice offers sight tests and fittings, you can book online.

\shot{13-appointments}{1.0}{Your appointments, with anything upcoming at
the top.}

\section{Making a booking}

\begin{enumerate}
  \item From your account, click \menu{Appointments}, then
        \btn{Book an appointment}.
  \item Choose a date and what the appointment is for.
  \item Click any available time.
\end{enumerate}

\shot{14-book-appointment}{1.0}{Choosing a time. Times already taken are
shown crossed out rather than hidden, so you can see how busy a day is.}

\section{Types of appointment}

\begin{longtable}{@{}p{0.22\textwidth}p{0.72\textwidth}@{}}
\toprule
\textbf{Type} & \textbf{What happens} \\
\midrule
\endhead
Eye test & A full sight test with an optometrist. \\
Fitting & Adjusting new glasses so they sit correctly. \\
Collection & Picking up finished glasses. \\
Adjustment & Straightening or tightening an existing pair. \\
Follow-up & A check after a change in prescription. \\
\bottomrule
\end{longtable}

\section{Opening hours}

The practice is open \textbf{10:00 to 18:00, Monday to Saturday}, and
closed on Sundays. Appointments can be booked up to 60 days ahead.

\section{Cancelling}

Click \btn{Cancel} beside the appointment. Please give as much notice as
you can --- the slot goes straight back into the diary for somebody
else.

% ---------------------------------------------------------------------
\chapter{Your data and your privacy}

Because VisionCart is an optical practice, it holds health information
about you. You have rights over that information, and the shop gives you
buttons rather than a form to fill in.

\shot{15-your-data}{1.0}{Your data, at \menu{Account} $\rightarrow$
\menu{Your data}.}

\section{Getting a copy of everything}

Click \btn{Download my data (JSON)}. You get a single file containing your
account details, addresses, orders, prescriptions and appointments.

\section{Asking for a change}

\shot{09-data-request}{0.92}{The request form. You do not need to be
signed in --- the right to ask does not depend on being able to get into
your account.}

You can ask the shop to:

\begin{itemize}
  \item \textbf{Correct} something that is wrong.
  \item \textbf{Send you a copy} of your data.
  \item \textbf{Restrict} how your data is used.
  \item \textbf{Erase} you from the records.
\end{itemize}

The shop replies within 30 days, and may check who you are first.

\section{What ``erase'' really means}

\begin{warnbox}
Erasure removes everything that identifies you --- your name, email,
telephone number, date of birth and addresses.

It does \textbf{not} delete your prescriptions or your past orders.
A prescription is a medical record and an order is a financial one, and
both have to be kept for a period the law sets. What is kept can no
longer be connected to you by name.

Erasure cannot be undone, and it closes your account.
\end{warnbox}

% =====================================================================
\part{Running the shop}
% =====================================================================

\chapter{Signing in to the back office}

The back office is where the shop is run. Everything in Part II happens
there.

\shot{06-signin}{0.92}{The sign-in page --- the same one customers use.
What you can see afterwards depends on your account.}

Go to your shop's address followed by \texttt{/admin}, or sign in
normally and you will be taken there.

\section{The four kinds of account}

\begin{longtable}{@{}p{0.20\textwidth}p{0.74\textwidth}@{}}
\toprule
\textbf{Role} & \textbf{Can do} \\
\midrule
\endhead
\textbf{Customer} &
Shop, manage their own account. No access to the back office at all. \\[5pt]
\textbf{Staff} &
Orders, patients, the diary, the catalogue, media, offers, delivery
charges, import and export, data requests. \\[5pt]
\textbf{Optician} &
Everything Staff can do, \emph{plus} approving and rejecting
prescriptions. \\[5pt]
\textbf{Administrator} &
Everything, plus shop settings, the audit trail and erasing customer
data. \\
\bottomrule
\end{longtable}

\begin{notebox}
If you cannot see a menu item described in this manual, your account
does not have that role. That is deliberate --- ask an administrator
rather than sharing somebody else's sign-in.
\end{notebox}

\begin{warnbox}
Never share an account. Every change is recorded against whoever was
signed in, and a shared login makes that record worthless exactly when
you need it.
\end{warnbox}

% ---------------------------------------------------------------------
\chapter{The dashboard}

The dashboard is the first thing you see, and it is designed to answer
one question: \emph{what needs attention today?}

\shot{20-dashboard}{1.0}{The dashboard. Tiles shaded amber are the ones
waiting on somebody.}

\section{The tiles}

\begin{longtable}{@{}p{0.30\textwidth}p{0.64\textwidth}@{}}
\toprule
\textbf{Tile} & \textbf{Meaning} \\
\midrule
\endhead
Orders today & Orders placed since midnight. \\
Paid, last 30 days & Money actually received. \\
\textbf{Awaiting payment} & Orders placed but not yet paid for. \\
In the lab & Orders with the lens laboratory now. \\
Patient files & How many people the practice has on record. \\
Live frames & Frames currently on sale. \\
\textbf{Prescriptions to check} & Waiting for an optician. \\
\textbf{Low stock lines} & Colourways nearly sold out. \\
\bottomrule
\end{longtable}

The amber tiles are work. The others are information.

\section{Prescriptions waiting for an optician}

Beneath the tiles is the queue of prescriptions waiting to be checked,
oldest first. This is the most important list in the shop --- nothing
can be made until an optician has been through it.

% ---------------------------------------------------------------------
\chapter{Orders, day to day}

\shot{21-orders}{1.0}{The order list. Payment and progress are shown
separately, because an order can be paid but not yet made, or made but
not yet paid.}

\section{Finding an order}

Use the search box for an order number, an email address or a telephone
number. The three drop-downs narrow the list by progress, payment or lab
stage.

\section{The life of an order}

\begin{longtable}{@{}p{0.20\textwidth}p{0.74\textwidth}@{}}
\toprule
\textbf{Stage} & \textbf{Meaning} \\
\midrule
\endhead
Pending & Placed, not yet paid. \\
Paid & Payment received. \\
In lab & With the laboratory, being made. \\
Ready & Made, waiting to be sent or collected. \\
Shipped & On its way. \\
Delivered & Complete. \\
Cancelled & Stopped. Stock goes back on the shelf. \\
Refunded & Money returned. \\
\bottomrule
\end{longtable}

\section{Taking a payment by hand}

For cash on delivery or a bank transfer, open the order and click
\btn{Mark as paid}, adding a reference such as the transfer number.

\begin{notebox}
Marking an order paid by hand does exactly what an automatic card
payment does --- same record, same email to the customer. There is no
second, weaker path through the shop.
\end{notebox}

\section{Sending an order out}

Open the order, click \btn{Mark as shipped}, choose the courier and enter the
tracking number. The customer is emailed automatically.

\section{Cancelling an order}

Click \btn{Cancel} and give a reason.

\begin{warnbox}
Cancelling returns the frames to stock. If the lenses have already been
made to a prescription, they cannot be sold to anybody else --- check
with the lab before cancelling anything that has passed \emph{In lab}.
\end{warnbox}

% ---------------------------------------------------------------------
\chapter{Checking prescriptions}

\begin{notebox}
This chapter is for \textbf{opticians}. Other staff can see
prescriptions but cannot approve them.
\end{notebox}

Every prescription typed in by a customer or imported from a spreadsheet
arrives as \textbf{pending verification}. Nothing is made until an
optician has looked at it.

\section{Working through the queue}

From the dashboard, click any patient in
\menu{Prescriptions waiting for an optician}, or go to \menu{Patients}
and open a file.

Check the values against whatever the customer supplied, then:

\begin{itemize}
  \item \btn{Verify} --- the order moves on to the lab, and the
        customer is emailed.
  \item \btn{Send \& reject} --- give a reason. The customer is emailed the
        reason, so write it as something they can act on: ``the axis is
        missing for the right eye'' rather than ``invalid''.
\end{itemize}

\section{Why a prescription can never be edited}

\begin{warnbox}
Once a prescription has been used on an order it is locked. If something
needs to change, you create a \emph{new version}; the old one stays
exactly as it was.

This is not awkwardness. If a customer ever reports a problem with their
glasses, the practice must be able to show precisely what was made and
what it was made from. A record that can be quietly edited afterwards
proves nothing.
\end{warnbox}

\section{Changes the shop will not accept}

Two rules are enforced everywhere --- the customer's form, your screen,
and spreadsheet imports alike:

\begin{itemize}
  \item Every value must sit on a 0.25 step.
  \item A cylinder must have an axis.
\end{itemize}

Both exist because a lab cannot make anything else.

% ---------------------------------------------------------------------
\chapter{Patient files}

\shot{22-patients}{1.0}{The patient list. Everyone who has ever ordered
has a file, including guests.}

Each file holds the person's contact details, every prescription they
have had, their orders, appointments and any documents.

\section{Finding somebody}

Search by name, file number, email or telephone. File numbers look like
\texttt{P-000042}.

\section{Adding a prescription}

Open the file and click \btn{Add prescription}. This creates a new
version --- it never overwrites what is there.

\begin{warnbox}
Patient files contain health information. Do not leave one open on a
screen the public can see, and never email a screenshot of one.
\end{warnbox}

% ---------------------------------------------------------------------
\chapter{The appointment diary}

\shot{23-diary}{1.0}{The diary. Move through the calendar with
\btn{Earlier}, \btn{Today} and \btn{Later}, and change how many days are
shown.}

\section{During the day}

Each booking has three buttons:

\begin{itemize}
  \item \btn{Seen} --- the patient attended.
  \item \btn{No show} --- they did not.
  \item \btn{Cancel} --- the slot goes back into the diary.
\end{itemize}

\section{Booking somebody in over the telephone}

Use the booking form on the diary page: choose the patient, the date and
time, the type and the length.

The diary will not let you double-book. If two bookings would overlap
for the same clinician, the second is refused with an explanation.

\begin{tipbox}
\btn{Seen} only becomes available once the appointment time has passed.
If you cannot click it, check you are not looking at tomorrow.
\end{tipbox}

% ---------------------------------------------------------------------
\chapter{The frame catalogue}

\shot{24-frames}{1.0}{The frame list, with stock levels.}

\section{Adding a frame}

Click \btn{New frame}.

\shot{25-frame-form}{1.0}{The frame form.}

\begin{longtable}{@{}p{0.26\textwidth}p{0.68\textwidth}@{}}
\toprule
\textbf{Field} & \textbf{Notes} \\
\midrule
\endhead
Name and SKU & The SKU is your own product code and must be unique. \\
Brand & Created automatically if it is new. \\
Price & In rupees, as you would write it on a price list. \\
Shape, material, rim & Used by the shop's filters. \\
Measurements & Lens width, bridge and arm length in millimetres. \\
Status & \emph{Draft} is invisible to customers; \emph{Active} is on
sale; \emph{Archived} is withdrawn but kept for old orders. \\
\bottomrule
\end{longtable}

\section{Colourways}

Each frame can have several colours. Each has its own product code,
stock level and photographs, so you can sell out of black without
affecting tortoiseshell.

\begin{notebox}
Prices are typed the way you would say them --- \texttt{8500} means
Rs.\,8,500. The shop converts and stores them in a way that cannot drift
by a rupee no matter how many discounts are applied.
\end{notebox}

% ---------------------------------------------------------------------
\chapter{Lens options and prices}

\shot{26-lenses}{1.0}{Lens options, grouped by type, thickness, coating
and tint.}

Click \btn{Add an option} to create one. Each option has a code, a
name, a price and a position in its group.
Turning one off hides it from customers immediately but leaves past
orders untouched.

% ---------------------------------------------------------------------
\chapter{Photographs and try-on artwork}

\shot{27-media}{1.0}{The media library.}

\section{Uploading}

Drag a whole shoot onto the drop area, or click to choose files. Each
photograph is:

\begin{itemize}
  \item turned the right way up automatically,
  \item reduced to a sensible size,
  \item converted to a modern format that loads quickly,
  \item given a thumbnail.
\end{itemize}

Files upload one at a time, so a single bad file is reported by name
while the rest of the shoot goes through.

\section{Try-on artwork}

Artwork for the virtual try-on must be a PNG with a transparent
background. Tick \field{Keep transparency} before uploading, or the
frame will appear with a white box around it.

\begin{warnbox}
An image attached to a colourway cannot be deleted --- the shop refuses,
rather than leaving a product page with a broken picture. Detach it
first.
\end{warnbox}

% ---------------------------------------------------------------------
\chapter{Offers and discount codes}

\shot{28-promotions}{1.0}{Promotions.}

Click \btn{New deal} to create one.

\section{Kinds of offer}

\begin{itemize}
  \item \textbf{Percentage off} the basket.
  \item \textbf{Fixed amount off}.
  \item \textbf{Free delivery}.
  \item \textbf{Buy one, get one} --- the cheaper pair is the free one.
\end{itemize}

\section{The settings that matter}

\begin{longtable}{@{}p{0.26\textwidth}p{0.68\textwidth}@{}}
\toprule
\textbf{Setting} & \textbf{What it does} \\
\midrule
\endhead
Code & Leave blank for an offer that applies automatically. \\
Minimum spend & Below this, the offer does not apply --- and the
customer is told why. \\
Maximum discount & Caps how much can come off. \\
Can be combined & Whether it may stack with other offers. \\
Priority & Which offer wins when several could apply. \\
Starts / ends & Leave blank to run indefinitely. \\
\bottomrule
\end{longtable}

\begin{warnbox}
An offer with no minimum spend, no cap and no end date is how shops lose
money. Set at least an end date.
\end{warnbox}

% ---------------------------------------------------------------------
\chapter{Delivery charges}

\shot{29-delivery}{1.0}{Delivery rates.}

Click \btn{Add a delivery rate} to create one. Each rate has a name
the customer sees, the area it applies to, a price
and an estimated number of days. Set the price to zero for free
delivery.

% ---------------------------------------------------------------------
\chapter{Spreadsheets: importing and exporting}

\shot{30-import-export}{1.0}{Import and export.}

\section{Exporting}

Four exports are available --- frames and stock, patients, prescriptions
and orders. Each downloads as a spreadsheet file you can open in Excel.

\begin{warnbox}
Patient and prescription exports contain health information. Every
download is recorded. Do not email these files or leave them in a
shared folder.
\end{warnbox}

\section{Importing}

\begin{enumerate}
  \item Export the matching file first --- it already has the right
        columns.
  \item Edit it in Excel.
  \item Choose the file, then click \btn{Check the file}.
  \item Read the report. Nothing has been written yet.
  \item If it looks right, click \btn{Import for real}.
\end{enumerate}

\begin{tipbox}
\btn{Check the file} lists problems by the line number you see in Excel,
so you can jump straight to the row and fix it.
\end{tipbox}

\section{How rows are matched}

Frames are matched on the colourway's product code, patients on their
file number. A row that matches updates the existing record; a row that
does not creates a new one.

\begin{notebox}
A bad row does not stop the good ones. If ten rows are fine and one is
wrong, the ten are imported and the one is reported.
\end{notebox}

% ---------------------------------------------------------------------
\chapter{Customer data requests}

\shot{31-data-requests}{1.0}{The data-request queue, oldest first ---
there is a legal clock running on each one.}

\section{Working through a request}

Open it, confirm who the person is, do what they asked, then set the
status to \emph{Completed} and write a note saying what you did and how
you confirmed their identity.

\begin{warnbox}
\textbf{Erasure is permanent and administrator-only.} It removes
everything identifying the customer across their file, orders and
addresses, and closes their account. Prescriptions and order totals are
kept, but can no longer be traced to a person.

The shop asks you to type \texttt{ERASE} to confirm, and refuses
entirely while the customer has an order still on its way --- the
courier needs somewhere to deliver to.
\end{warnbox}

% =====================================================================
\part{Looking after the shop}
% =====================================================================

\chapter{Settings}

\shot{33-settings}{1.0}{Shop settings. Administrators only.}

Here you set the shop's name, currency, tax rate, contact details, which
payment methods are offered and the text customers see for bank
transfers.

\begin{warnbox}
Changing the tax rate affects every new order immediately. It does not
change orders already placed --- those keep the figures they were sold
at, which is exactly what an accountant needs.
\end{warnbox}

% ---------------------------------------------------------------------
\chapter{The audit trail}

\shot{32-audit}{1.0}{The audit trail. Administrators only.}

Every change to a patient record, a price or an order is recorded: who
did it, what they did, when, and from which computer.

\section{What it is for}

\begin{itemize}
  \item Answering ``who changed this price?''
  \item Showing a regulator that patient records are controlled.
  \item Investigating a mistake without guesswork.
\end{itemize}

\begin{notebox}
The trail deliberately records \emph{that} a prescription was changed,
never the values themselves. Clinical detail lives on the patient file,
where access is controlled --- not in a log that staff browse casually.
\end{notebox}

% ---------------------------------------------------------------------
\chapter{Everyday care}

A short list of things worth doing regularly.

\section{Each day}

\begin{itemize}
  \item Clear the \textbf{Prescriptions to check} queue. Nothing gets
        made until you do.
  \item Look at \textbf{Awaiting payment} for orders that have stalled.
  \item Check the diary for tomorrow.
\end{itemize}

\section{Each week}

\begin{itemize}
  \item Review \textbf{Low stock lines} and reorder.
  \item Check the data-request queue --- the legal clock does not stop
        for a busy week.
  \item Glance at the audit trail for anything unexpected.
\end{itemize}

\section{Each month}

\begin{itemize}
  \item Export your orders for your accountant.
  \item Review which offers are still running.
  \item Confirm your backups actually exist.
\end{itemize}

\begin{warnbox}
\textbf{A database backup on its own is not a backup.} Your product
photographs live alongside it, not inside it. Restoring one without the
other gives you a catalogue of broken pictures. Back up both, and check
once in a while that you can actually restore them.
\end{warnbox}

\section{Is the shop healthy?}

Two addresses answer that without needing anyone technical:

\begin{itemize}
  \item \texttt{/health/live} --- is the shop switched on?
  \item \texttt{/health/ready} --- can it actually take an order?
\end{itemize}

The second is the one that matters. A shop can be switched on and still
be unable to reach its records, and from the outside those look
identical.

% =====================================================================
\appendix
% =====================================================================

\chapter{Glossary}

\begin{longtable}{@{}p{0.26\textwidth}p{0.68\textwidth}@{}}
\toprule
\textbf{Term} & \textbf{Meaning} \\
\midrule
\endhead
\textbf{ADD} & Extra reading power, on bifocal and varifocal prescriptions. \\
\textbf{Anti-reflective} & A coating that cuts glare from lights and screens. \\
\textbf{Astigmatism} & An eye shaped slightly more like a rugby ball than a football. Corrected by the cylinder and axis. \\
\textbf{Axis} & An angle from 1 to 180 telling the lab which way round to grind a cylinder. \\
\textbf{Bifocal} & Two prescriptions in one lens with a visible line. \\
\textbf{Bridge} & The part over your nose. Measured in millimetres. \\
\textbf{Colourway} & One colour of a frame. Each has its own code, stock and pictures. \\
\textbf{CYL} & Cylinder. The astigmatism correction. \\
\textbf{Diopter} & The unit lens strength is measured in. Always in steps of 0.25. \\
\textbf{Index} & How efficiently a lens bends light. Higher means thinner. \\
\textbf{OD} & Your right eye. \\
\textbf{OS} & Your left eye. \\
\textbf{Patient file} & The practice's record of one person. \\
\textbf{PD} & Pupillary distance --- between the centres of your pupils, in millimetres. \\
\textbf{Photochromic} & Lenses that darken in sunlight. \\
\textbf{Single vision} & One prescription across the whole lens. \\
\textbf{SKU} & A product code. \\
\textbf{SPH} & Sphere. How short- or long-sighted you are. \\
\textbf{Varifocal} & Distance, middle and reading in one lens, with no line. \\
\bottomrule
\end{longtable}

% ---------------------------------------------------------------------
\chapter{If something goes wrong}

\section*{For customers}

\begin{longtable}{@{}p{0.34\textwidth}p{0.60\textwidth}@{}}
\toprule
\textbf{Problem} & \textbf{Try this} \\
\midrule
\endhead
The camera will not start &
Allow camera access from the padlock icon in your browser's address bar,
and close any video-call software. \\[5pt]
The frame sits crookedly &
Click \btn{Adjust eye points} and mark the centre of each pupil. \\[5pt]
My discount code is refused &
The shop says why. Usually the code has expired or the basket is below
the minimum. \\[5pt]
I cannot enter my prescription &
Values step by 0.25 and a cylinder needs an axis. Check your
prescription again. \\[5pt]
No confirmation email &
Check your spam folder, then contact the shop with your order number. \\[5pt]
I have forgotten my password &
\btn{Forgotten your password?} on the sign-in page. \\
\bottomrule
\end{longtable}

\section*{For staff}

\begin{longtable}{@{}p{0.34\textwidth}p{0.60\textwidth}@{}}
\toprule
\textbf{Problem} & \textbf{Try this} \\
\midrule
\endhead
A menu item is missing &
Your account does not have that role. Ask an administrator. \\[5pt]
``That slot has just been taken'' &
Somebody booked it while you were typing. Choose another time. \\[5pt]
An import failed &
Read the line numbers in the report and fix those rows in Excel. \\[5pt]
An image will not delete &
It is attached to a colourway. Detach it first. \\[5pt]
Lenses cannot be marked ready &
The prescription has not been approved by an optician yet. \\[5pt]
Too many sign-in attempts &
Sign-in is limited to protect accounts. Wait five minutes. \\
\bottomrule
\end{longtable}

\vspace{1.4cm}
\begin{center}
\color{vcgrey}\small
\rule{0.35\textwidth}{0.4pt}\\[10pt]
\emph{VisionCart Optical --- User Manual}\\
August 2026
\end{center}

\end{document}
